\section{Zoo FHE Architecture}

\label{sec:architecture}

\subsection{Dual-Scheme Design}

Zoo FHE uses both \TFHE{} and \CKKS{} within a unified framework:

\begin{figure}[t]
\centering
\begin{tikzpicture}[
    node distance=1.5cm,
    box/.style={rectangle, draw, minimum width=3cm, minimum height=0.9cm, align=center, fill=blue!10},
    arrow/.style={->, thick}
]
    \node[box] (input) {Encrypted Input\\(\CKKS{} packed)};
    \node[box, right=2cm of input] (conv) {Convolution\\(\CKKS{} SIMD)};
    \node[box, right=2cm of conv] (relu) {ReLU\\(\TFHE{} bootstrap)};
    \node[box, below=of relu] (pool) {Pooling\\(\CKKS{} rotation)};
    \node[box, left=2cm of pool] (fc) {Dense Layer\\(\CKKS{} matmul)};
    \node[box, left=2cm of fc] (output) {Encrypted Output\\(Threshold Dec)};

    \draw[arrow] (input) -- (conv);
    \draw[arrow] (conv) -- (relu);
    \draw[arrow] (relu) -- (pool);
    \draw[arrow] (pool) -- (fc);
    \draw[arrow] (fc) -- (output);
\end{tikzpicture}
\caption{Dual-scheme neural network inference. Linear layers use \CKKS{} SIMD; nonlinear activations use \TFHE{} bootstrapping.}
\label{fig:dual-scheme}
\end{figure}

The key insight is that neural networks alternate between linear operations (matrix multiplications, convolutions) and nonlinear activations (ReLU, sigmoid). \CKKS{} efficiently handles linear operations via SIMD packing, while \TFHE{} handles nonlinear operations via programmable bootstrapping.

\subsection{Scheme Switching}

Transitioning between \CKKS{} and \TFHE{} ciphertexts requires a bridging protocol:

\begin{algorithm}[t]
\caption{CKKS-to-TFHE Bridge}
\label{alg:bridge}
\begin{algorithmic}[1]
\STATE \textbf{Input}: \CKKS{} ciphertext $c_{\text{ckks}}$ encoding value $x \in [-1, 1]$
\STATE \textbf{Output}: \TFHE{} ciphertext $c_{\text{tfhe}}$ encoding $\lfloor x \cdot 2^p \rfloor$
\STATE Scale: $c' \leftarrow c_{\text{ckks}} \cdot 2^p$
\STATE Extract coefficients from $c'$
\STATE Construct LWE ciphertext from coefficients
\STATE $c_{\text{tfhe}} \leftarrow \text{KeySwitch}(c_{\text{lwe}}, \text{ksk}_{\text{ckks} \to \text{tfhe}})$
\RETURN $c_{\text{tfhe}}$
\end{algorithmic}
\end{algorithm}

The key-switching key $\text{ksk}_{\text{ckks} \to \text{tfhe}}$ is generated during setup and stored encrypted on-chain.

\subsection{Precompile Interface}

The Zoo FHE precompile at \texttt{0x0700} exposes:

\begin{lstlisting}[basicstyle=\ttfamily\small,caption=Zoo FHE Precompile Interface]
interface IFHE {
    // CKKS operations (packed vectors)
    function ckksAdd(bytes ct1, bytes ct2)
        returns (bytes);
    function ckksMul(bytes ct1, bytes ct2)
        returns (bytes);
    function ckksRotate(bytes ct, int steps)
        returns (bytes);
    function ckksMatMul(bytes ct, bytes ptMatrix)
        returns (bytes);

    // TFHE operations (scalar integers)
    function tfheAdd(uint256 ct1, uint256 ct2)
        returns (uint256);
    function tfheLt(uint256 ct1, uint256 ct2)
        returns (uint256);
    function tfheMax(uint256 ct1, uint256 ct2)
        returns (uint256);

    // Scheme switching
    function ckksToTfhe(bytes ckksCt)
        returns (uint256);
    function tfheToCkks(uint256 tfheCt)
        returns (bytes);

    // Threshold decryption
    function requestDecrypt(uint256 ct)
        returns (bytes32 requestId);
}
\end{lstlisting}

